\begin{proposition}[B.4.12 (Proved): a proper quadric--parabolic surface smoothing]
Let $S=S_8(a,\kappa)$ be as in B.4.11, with elliptic axis $E$
and normal line $N$ of degree $-8$.
There is an embedding $E\hookrightarrow Q=\PP^1\times\PP^1$
as a smooth divisor of bidegree $(2,2)$ with
$N_{E/Q}\simeq N^{-1}$.
There is a proper flat map from a smooth complex threefold
\[
 \varpi:\mathcal V\longrightarrow\Delta
\]
with smooth noncentral fibres and reduced simple normal-crossings
central fibre exactly $S\cup_E Q$.  Its two normal lines are
\[
 N_{S/\mathcal V}=\OO_S(-E),\qquad
 N_{Q/\mathcal V}=K_Q=\OO_Q(-2,-2).
\]
The family is obtained by globalizing a specified cone smoothing
and blowing up its vertex; no finite base change is needed.

For the elliptic-quartic configuration of B.4.10, the three
pairwise normal-line products, including the triple-curve factors,
can consequently all be made trivial using
\[
 A=\operatorname{Bl}_E\PP^3,\qquad
 B=S\times\PP^1,\qquad C=\mathcal V.
\]
Here $C$ is a proper family over an open disc and is not a compact
threefold.  This assertion supplies the indicated surface smoothing
and normal-line products only.  It does not construct a compact
three-component central fibre, a compatible global smoothing of
such a fibre, or a comparison with the prescribed punctured
$S^6$ family.
\end{proposition}
\begin{proof}
Put $L=N^{-1}$, of degree eight.  Multiplication by two on the
elliptic Picard variety is surjective, so choose a degree-four
line bundle $M$ with $M^{\otimes2}=L$.
Choose degree-two line bundles $M_1,M_2$ with
$M_1\otimes M_2=M$ and $M_1\not\simeq M_2$.
Each complete linear system is a basepoint-free double cover of
$\PP^1$.  The product of the two maps is an embedding in $Q$.
Indeed two distinct points identified by both maps would give
$M_1\simeq\OO_E(p+q)\simeq M_2$.
A common ramification point would similarly give
$M_1\simeq\OO_E(2p)\simeq M_2$.
Thus the product is injective and immersive; properness makes it a
closed embedding.  Its image has bidegree $(2,2)$, and
\[
 \OO_Q(E)|_E=(M_1\otimes M_2)^{\otimes2}=L.
\]
This also allows us to start with the fixed quadric and elliptic
quartic in B.4.10: the normal line there was chosen to be the
inverse of this same $L$.

Use the anticanonical embedding of $Q$ in $\PP^8$ and write
\[
 R_Q=\bigoplus_{d\ge0}H^0(Q,\OO_Q(2d,2d)),\qquad
 Z=\operatorname{Spec}R_Q.
\]
The ring is generated in degree one and is normal.
One direct verification uses the bihomogeneous monomials:
every monomial of bidegree $(2d,2d)$ is a product of $d$
monomials of bidegree $(2,2)$, and the associated monoid is the
intersection of a rational polyhedral cone with its full lattice.
The punctured cone is the complement of the zero section in
$\OO_Q(-2,-2)$, hence is smooth.
Choose the section $e\in H^0(Q,\OO_Q(2,2))$ cutting out $E$.
The degree-one function $t=e$ on $Z$ defines a flat map
$Z\to\mathbb A^1$: its coordinate ring is torsion-free over
$\C[t]$, a principal ideal domain.
For every $d\ge0$, the sequence
\[
 0\longrightarrow\OO_Q(2d-2,2d-2)
 \xrightarrow{\ e\ }\OO_Q(2d,2d)
 \longrightarrow L^{\otimes d}\longrightarrow0
\]
has vanishing first cohomology in its left term.
For $d=0$ this follows from the Kunneth formula for
$\OO_Q(-2,-2)$; for $d\ge1$ it follows from the same formula
and nonnegative degrees.
Taking sections identifies
\[
 R_Q/(e)=\bigoplus_{d\ge0}H^0(E,L^{\otimes d}).
\]
The central surface is therefore exactly the affine cone over
$(E,L)$, including its scheme structure.
For $t\ne0$, projection from the cone vertex identifies the
fibre with $Q\setminus E$: over each point where $e$ is nonzero
there is exactly one radial coordinate on which $e$ has value $t$.
In particular this is an explicit smoothing of the isolated
singularity of B.4.11.
It is the affine version of Pinkham's cone construction
\citep[Remark~7.6(iii), p.~46; Section~9, pp.~63--65]{Pinkham1974Root}.
In degree eight the quadric component is distinguished from the
other smoothing component; no replacement of it by that other
component is made here.

Blowing up the vertex of $Z$ gives
$\operatorname{Tot}(\OO_Q(-2,-2))=\operatorname{Tot}(K_Q)$.
This can be checked on the blow-up charts given by ratios of the
degree-one homogeneous coordinates; these recover a point of
$Q$, with the remaining common factor the tautological fibre
coordinate $\xi$.
The map to the parameter line is precisely
\[
 t=\xi e.
\]
Its central divisor is the union of the zero section $Q$ and
the restriction $\operatorname{Tot}(K_Q|_E)$, each with
multiplicity one.  Near $E$ the equation is $t=\xi z$, where
$z$ is a local defining function of $E$ in $Q$.
Thus the total space is smooth and the central divisor has
simple normal crossings.

Now contract $E$ in the compact surface $S$ as in B.4.11.
The regular global-to-local deformation map in that proposition
lifts the specified analytic germ $Z\to\mathbb A^1$ to a proper
deformation $\mathcal Y\to\Delta$ of $Y=S/E$.
The lifted deformation agrees analytically with this cone
smoothing in a neighbourhood of its singular point.
Away from that neighbourhood the central surface is smooth;
openness of smoothness and properness allow a common shrinking
of $\Delta$ for which the family there is a submersion.
Blow up the single cone vertex in $\mathcal Y$.
The resulting map $\mathcal V\to\Delta$ remains proper.
Near the exceptional divisor its local model is the smooth
space just computed; elsewhere it is unchanged and smooth.
Its central strict transform is exactly $S$, because the
maximal-ideal blow-up of the degree-eight elliptic cone is its
given resolution $\operatorname{Tot}(N)$ near $E$.
The other component is $Q$, and their intersection is the original
marked elliptic curve.  The local equation $t=\xi z$ proves
flatness and reducedness near the intersection; at a smooth
point of either component $t$ is a local coordinate.
The noncentral fibres are the smooth compact fibres of
$\mathcal Y$.
This proves every asserted properness and smoothness property
for the surface family.

The exceptional zero section has normal $K_Q$ in its local
total-space model.  Since $S+Q$ is the principal divisor of $t$,
restriction to $S$ gives
\[
 \OO_{\mathcal V}(S)|_S
   =\OO_{\mathcal V}(-Q)|_S=\OO_S(-E).
\]
These are the two asserted normal lines.  They also show directly
that the two branch normals along $E$ multiply to the trivial line.

We record the precise extent of the three-piece matching.
Realize $Q$ as a smooth quadric in $\PP^3$ and $E$ as its
intersection with a second quadric.  Put $L=\OO_E(2)$.
Then $N_{E/\PP^3}=L\oplus L$, so the exceptional divisor
of $A=\operatorname{Bl}_E\PP^3$ is
$D_{AB}=E\times\PP^1$, with normal
$L\boxtimes\OO_{\PP^1}(-1)$.
The strict transform of $Q$ is still $Q$, because the centre is
Cartier on it, and its normal in $A$ is
\[
 N_{Q/A}=\OO_Q(2)\otimes\OO_Q(-E)=\OO_Q.
\]
It meets $D_{AB}$ in the constant section $E\times\{p\}$
corresponding to that member of the pencil of quadrics.
In $B=S\times\PP^1$, use the divisors
$D_{BA}=E\times\PP^1$ and $D_{BC}=S\times\{p\}$.
Their normals are $L^{-1}\boxtimes\OO$ and $\OO_S$.
In $C=\mathcal V$, use its central divisors $Q$ and $S$.
The three normal products with the triple contribution are then
\[
\begin{array}{c|c}
 \text{double surface}&
 N_{D/\text{first branch}}\otimes
 N_{D/\text{second branch}}\otimes\OO_D(\text{triple curve})\\ \hline
 E\times\PP^1&
 (L\boxtimes\OO(-1))\otimes(L^{-1}\boxtimes\OO)
       \otimes(\OO_E\boxtimes\OO(1))=\OO\\
 Q&\OO_Q\otimes K_Q\otimes\OO_Q(E)=\OO_Q\\
 S&\OO_S\otimes\OO_S(-E)\otimes\OO_S(E)=\OO_S.
\end{array}
\]
These are pairwise line-bundle identities.  They do not supply
a compact third component or a compatible global smoothing
section for a compact normal-crossings threefold.
In fact the third space constructed here maps surjectively and
properly to an open disc, so it cannot itself be compact.

For a useful topological check on the local model,
the affine smooth fibre $Q\setminus E$ has
\[
 H_1(Q\setminus E;\Z)=\Z/2,\qquad
 H_2(Q\setminus E;\Z)=\Z^3.
\]
Indeed the Thom sequence of the pair $(Q,Q\setminus E)$ has
intersection map $H_2(Q;\Z)\to H_0(E;\Z)$ given by
$(a,b)\mapsto2a+2b$.  Its cokernel is $\Z/2$ and its
kernel is $\Z$; the preceding term $H_1(E;\Z)=\Z^2$
injects into $H_2(Q\setminus E;\Z)$.
This exact extension of free groups splits abstractly.
The map $H_4(Q;\Z)\to H_2(E;\Z)$ is an isomorphism,
so $H_3(Q\setminus E;\Z)=0$, and its Euler number is four.
The Euler number is also obtained independently by projection
to either ruling of $Q$: away from the four branch values of
$E\to\PP^1$, the fibre is $\C^*$ and has Euler number zero;
at each branch value it is $\C$ and has Euler number one.
In particular no simply connected local Milnor fibre is being
assumed.  The compact noncentral surface has Euler number
$\chi(S)+\chi(Q)-2\chi(E)=8+4=12$ by the local
$\xi z=t$ retraction, not the Euler number of an $S^6$ fibre.
Exact polynomial, Hilbert-function, normal-product and Smith-form
checks appear in \texttt{code/root\_quadric\_smoothing.py}.
\end{proof}
